% ============================================================
% CHAPITRE 5
% ============================================================
\chapter{Sprint 4 : Événements et devis}

\section{Introduction}

Le Sprint 4 touche au cœur du réacteur financier du traiteur : la logistique des mariages, soutenances et fêtes d'entreprises. Vendre un plat à l'unité est une chose, mais orchestrer un buffet pour cent convives avec des contraintes de dates, de nappage et de tolérances alimentaires en est une autre. Ce volet démontre comment la plateforme s'émancipe d'une simple boutique en ligne pour devenir un véritable bureau d'organisation (gestion des réceptions et montages des devis).

% --------------------------------------------------
\section{Backlog du sprint}

\begin{table}[H]
    \centering
    \caption{Backlog du sprint 4}
    \label{tab:backlog_sprint4}
    \renewcommand{\arraystretch}{1.3}
    \begin{tabular}{|c|p{8cm}|c|}
        \hline
        \textbf{ID} & \textbf{User Story} & \textbf{Points} \\
        \hline
        US-20 & En tant que client, je souhaite soumettre une demande d'événement. & 8 \\
        \hline
        US-21 & En tant que client, je souhaite sélectionner des plats pour mon événement. & 5 \\
        \hline
        US-22 & En tant que client, je souhaite recevoir un devis pour mon événement. & 5 \\
        \hline
        US-23 & En tant qu'administrateur, je souhaite gérer les demandes d'événements. & 5 \\
        \hline
        US-24 & En tant qu'administrateur, je souhaite gérer les modèles d'événements. & 3 \\
        \hline
    \end{tabular}
\end{table}

% --------------------------------------------------
\section{Conception}

\subsection{Diagramme de cas d'utilisation}

\begin{figure}[H]
    \centering
    \includegraphics[width=0.85\textwidth]{uc_sprint4.png}
    \caption{Diagramme de cas d'utilisation du sprint 4}
    \label{fig:uc_sprint4}
\end{figure}

\subsection{Description textuelle des cas d'utilisation}

\subsubsection{Cas d'utilisation : Soumettre une demande d'événement}

\begin{table}[H]
    \centering
    \caption{Description textuelle --- Soumettre une demande d'événement}
    \label{tab:desc_demande_event}
    \renewcommand{\arraystretch}{1.3}
    \begin{tabular}{|p{4cm}|p{8.5cm}|}
        \hline
        \textbf{Cas d'utilisation} & Soumettre une demande d'événement \\
        \hline
        \textbf{Acteur principal} & Client \\
        \hline
        \textbf{Précondition} & Le client est authentifié \\
        \hline
        \textbf{Postcondition} & La demande est créée avec le statut PENDING \\
        \hline
        \textbf{Scénario principal} &
        \begin{enumerate}
            \item Le futur marié (ou l'organisateur) ouvre le module de planification
            \item Spécification de la nature de la cérémonie (Fiançailles, Séminaire)
            \item Paramétrage du nombre d'invités, de l'adresse de la salle et du jour J
            \item Saisie des exigences particulières (ex: "Aucune viande rouge")
            \item Sélection minutieuse de la composition du buffet
            \item L'algorithme dresse une enveloppe budgétaire brute
            \item Envoi du dossier complet à la direction
            \item Cloche de notification chez l'administrateur
        \end{enumerate} \\
        \hline
        \textbf{Scénario alternatif} &
        \begin{itemize}
            \item[3a.] Réservation dans le passé impossible : le calendrier grise ces jours
            \item[3b.] Surcharge de la capacité maximale du traiteur : jauge bloquante
        \end{itemize} \\
        \hline
    \end{tabular}
\end{table}

\subsubsection{Cas d'utilisation : Gérer les demandes d'événements (Administrateur)}

\begin{table}[H]
    \centering
    \caption{Description textuelle --- Gérer les demandes d'événements}
    \label{tab:desc_gerer_events}
    \renewcommand{\arraystretch}{1.3}
    \begin{tabular}{|p{4cm}|p{8.5cm}|}
        \hline
        \textbf{Cas d'utilisation} & Gérer les demandes d'événements \\
        \hline
        \textbf{Acteur principal} & Administrateur \\
        \hline
        \textbf{Précondition} & L'administrateur est authentifié \\
        \hline
        \textbf{Postcondition} & La demande est traitée (confirmée, annulée ou maintenue en attente) \\
        \hline
        \textbf{Scénario principal} &
        \begin{enumerate}
            \item Interception du dossier de réception par le pôle administratif
            \item Ouverture de la fiche détaillée du futur banquet
            \item Audit de la jauge d'invités et des spécificités demandées
            \item Surfacturation éventuelle (frais kilométriques, service à table)
            \item Cachet final de l'entreprise : "Accepté" ou "Désolé, indisponible"
            \item Le client voit la couleur de son badge changer dans son historique
        \end{enumerate} \\
        \hline
        \textbf{Scénario alternatif} &
        \begin{itemize}
            \item[5a.] Blocage du devis : expédition d'un courrier d'excuse argumenté
        \end{itemize} \\
        \hline
    \end{tabular}
\end{table}

\subsection{Diagrammes de séquence}

\subsubsection{Séquence : Soumission d'une demande d'événement}

\begin{figure}[H]
    \centering
    \includegraphics[width=0.9\textwidth]{seq_demande_event.png}
    \caption{Diagramme de séquence --- Soumission d'une demande d'événement}
    \label{fig:seq_demande_event}
\end{figure}

\subsubsection{Séquence : Traitement d'une demande par l'administrateur}

\begin{figure}[H]
    \centering
    \includegraphics[width=0.9\textwidth]{seq_traitement_event.png}
    \caption{Diagramme de séquence --- Traitement d'une demande d'événement}
    \label{fig:seq_traitement_event}
\end{figure}

\subsection{Diagramme de classes}

\begin{figure}[H]
    \centering
    \includegraphics[width=0.85\textwidth]{class_sprint4.png}
    \caption{Diagramme de classes du service d'événements}
    \label{fig:class_sprint4}
\end{figure}

Le service d'événements repose sur trois entités :

\textbf{Le modèle \texttt{Event}} contient :
\begin{itemize}
    \item \texttt{id} : identifiant unique (UUID)
    \item \texttt{userId} : identifiant du client demandeur
    \item \texttt{type} : type d'événement (\texttt{WEDDING}, \texttt{BIRTHDAY}, \texttt{CORPORATE}, \texttt{OTHER})
    \item \texttt{date} : date de l'événement
    \item \texttt{location} : lieu de l'événement
    \item \texttt{guestCount} : nombre de convives
    \item \texttt{notes} : notes complémentaires du client
    \item \texttt{status} : état de la demande (\texttt{PENDING}, \texttt{CONFIRMED}, \texttt{CANCELLED})
    \item \texttt{totalEstimate} : estimation du coût total (type décimal)
    \item \texttt{items} : collection des plats sélectionnés
\end{itemize}

\textbf{Le modèle \texttt{EventItem}} contient :
\begin{itemize}
    \item \texttt{id} : identifiant unique
    \item \texttt{platId} : identifiant du plat sélectionné
    \item \texttt{platName} : nom du plat (dénormalisé)
    \item \texttt{platPrice} : prix unitaire du plat
    \item \texttt{quantity} : quantité souhaitée
    \item \texttt{eventId} : clé étrangère vers \texttt{Event}
\end{itemize}

\textbf{Le modèle \texttt{EventTemplate}} contient :
\begin{itemize}
    \item \texttt{id} : identifiant unique
    \item \texttt{name} : nom du modèle (ex. : <<~Mariage classique~>>)
    \item \texttt{description} : description du modèle
    \item \texttt{type} : type d'événement associé
    \item \texttt{defaultGuestCount} : nombre de convives par défaut
    \item \texttt{suggestedItems} : liste de plats suggérés (JSON)
\end{itemize}

\subsection{Diagramme d'états-transitions}

Le diagramme d'états-transitions (figure~\ref{fig:state_sprint4}) modélise le cycle de vie d'une demande d'événement, depuis la soumission par le client jusqu'à la clôture de la cérémonie.

\begin{figure}[H]
    \centering
    \includegraphics[width=0.85\textwidth]{state_sprint4_event.png}
    \caption{Diagramme d'états-transitions --- Cycle de vie d'un événement}
    \label{fig:state_sprint4}
\end{figure}

% --------------------------------------------------
\section{Réalisation}

\subsection{Points d'entrée de l'API}

\begin{table}[H]
    \centering
    \caption{Endpoints REST du service d'événements}
    \label{tab:endpoints_event}
    \renewcommand{\arraystretch}{1.2}
    \begin{tabular}{|l|p{5cm}|p{4.5cm}|}
        \hline
        \textbf{Méthode} & \textbf{Route} & \textbf{Description} \\
        \hline
        POST & /api/events & Soumission d'une demande \\
        \hline
        GET & /api/events & Événements du client \\
        \hline
        GET & /api/events/:id & Détails d'un événement \\
        \hline
        GET & /api/events/templates & Liste des modèles \\
        \hline
        GET & /api/events/admin/all & Tous les événements (administrateur) \\
        \hline
        PUT & /api/events/admin/:id/status & Mise à jour du statut (administrateur) \\
        \hline
        POST & /api/events/admin/templates & Création d'un modèle (administrateur) \\
        \hline
        PUT & /api/events/admin/templates/:id & Modification d'un modèle (administrateur) \\
        \hline
        DELETE & /api/events/admin/templates/:id & Suppression d'un modèle (administrateur) \\
        \hline
    \end{tabular}
\end{table}

% Schéma de la base de données retiré : le diagramme de classes (figure~\ref{fig:class_sprint4}) suffit à modéliser la structure des données.



\subsection{Calcul du devis}

L'estimation du coût d'un événement suit la formule :

$$
\text{Total} = \sum_{i=1}^{n} (prix_i \times quantité_i) + \text{frais de service}
$$

où $n$ est le nombre de plats sélectionnés, $prix_i$ est le prix unitaire du $i$-ème plat et $quantité_i$ la quantité demandée. Les frais de service peuvent être ajustés par l'administrateur lors de la validation du devis.

\subsection{Interfaces réalisées}

\subsubsection{Formulaire de demande d'événement}

\begin{figure}[H]
    \centering
    \includegraphics[width=0.85\textwidth]{ui_demande_event.png}
    \caption{Interface de soumission d'une demande d'événement}
    \label{fig:ui_demande_event}
\end{figure}

Le parcours de réservation accompagne le futur marié (ou l'organisateur) étape par étape : d'abord le choix du type de célébration, puis le renseignement du lieu, de la date et de la jauge d'invités. Vient ensuite la composition du buffet depuis le catalogue Gala. Un encadré chiffré synthétise l'enveloppe budgétaire avant l'envoi définitif du dossier.

\subsubsection{Sélection de plats pour un événement}

\begin{figure}[H]
    \centering
    \includegraphics[width=0.85\textwidth]{ui_selection_plats_event.png}
    \caption{Interface de sélection de plats pour un événement}
    \label{fig:ui_selection_plats_event}
\end{figure}

\subsubsection{Suivi des demandes (Client)}

\begin{figure}[H]
    \centering
    \includegraphics[width=0.85\textwidth]{ui_suivi_events.png}
    \caption{Interface de suivi des demandes d'événements}
    \label{fig:ui_suivi_events}
\end{figure}

Le carnet privé du client empile toutes les réceptions soumises depuis son compte. Chaque ligne arbore une pastille : jaune si le devis est en cours de révision, verte si la direction a validé le menu, rouge si l'événement a été décliné.

\subsubsection{Gestion des événements (Administrateur)}

\begin{figure}[H]
    \centering
    \includegraphics[width=0.9\textwidth]{ui_admin_events.png}
    \caption{Interface d'administration des événements}
    \label{fig:ui_admin_events}
\end{figure}

Sur l'écran de M. Daoud, l'ensemble des dossiers de réception apparaît avec des filtres par nature de cérémonie (mariage, séminaire) et par état. Un clic ouvre la fiche détaillée du banquet, où le gérant peut retailler le chiffrage et entériner ou refuser la prestation.

% --------------------------------------------------
\section{Tests fonctionnels}

\begin{table}[H]
    \centering
    \caption{Tests fonctionnels --- Sprint 4}
    \label{tab:tests_sprint4}
    \renewcommand{\arraystretch}{1.3}
    \begin{tabular}{|c|p{3cm}|p{3.5cm}|p{2.5cm}|}
        \hline
        \textbf{N°} & \textbf{Cas de test} & \textbf{Comportement attendu} & \textbf{Résultat} \\
        \hline
        1 & Soumission d'une demande complète & Demande créée avec statut PENDING & Réussi \\
        \hline
        2 & Soumission avec date passée & Message d'erreur de validation & Réussi \\
        \hline
        3 & Sélection de plats pour un événement & Plats ajoutés et coût estimé calculé & Réussi \\
        \hline
        4 & Consultation des événements du client & Liste des demandes avec statuts corrects & Réussi \\
        \hline
        5 & Mise à jour du statut par l'administrateur & Statut modifié en base & Réussi \\
        \hline
        6 & Création d'un modèle d'événement & Modèle créé et disponible à la sélection & Réussi \\
        \hline
        7 & Suppression d'un modèle & Modèle supprimé du système & Réussi \\
        \hline
        8 & Accès aux événements sans authentification & Redirection vers la connexion & Réussi \\
        \hline
    \end{tabular}
\end{table}

% --------------------------------------------------
\section{Conclusion}

Avec ce Sprint 4 achevé, le système prouve sa vraie valeur ajoutée pour un métier de bouche supérieur. Fini le carnet de commandes chaotique ou les brouillons égarés : tout internaute peut concevoir l'architecture de son dîner de gala directement depuis son canapé, tandis que la direction ajuste l'offre monétaire sans quitter son tableau de bord. Le socle étant totalement posé, le Sprint final va pouvoir se concentrer sur la cerise technologique de ce PFE : l'arrivée de l'assistant robotique (IA).
